\chapter{Qt}
\begin{minted}{cpp}

//UI界面的界面布局可以添加Widget窗口后，再选中mainwindow,点击水平或垂直布局

 //::在C++表示作用域解析运算符,用于明确指定某个标识符属于哪个作用域
/* void Widget::initPortSettings()
{

} 
不是在调用函数，而是在定义函数
*/

int main(int argc, char *argv[])//必须使用两个参数的main函数，通过参数创建Qt应用程序
{
    QApplication a(argc, argv);//应用程序实例对象，有且只能一个
    Widget w;//创建窗口对象w
    /*创建类的对象有两种方法
     1.如像上面的方式创建：在栈内存上创建对象，对象可以使用.(成员运算符)直接访问成员和方法
    是一种静态联编，为xxx创建了地址空间
     2.使用指针创建：在堆内存上创建对象
    Widget* window;
    window=new Widget;
    ->等价于（*window）.;
    */
    w.show();//显示窗口
    return a.exec();//Qt进入消息循环
}

/*
对于创建对象实例两种方法的比喻和理解：

*/

//另一个.cpp文件里的代码，这里学习下构造函数和初始化列表的知识
Widget::Widget(QWidget *parent)
    : QWidget(parent)//初始化函数列表1
    , ui(new Ui::Widget)//初始化列表2
{
    ui->setupUi(this);
    //设置窗口标题
    this->setWindowTitle("雷达参数配置及CLI命令设置界面");
    //设置窗口图标
    // this->setWindowIcon(QIcon());
    initPortSettings();

    //创建对象并实例化：监听的服务器对象
    m_s= new QTcpServer(this);
    //启动监听服务；；；；这里注意：使用传到槽就不要使用connect，如果使用connect就不要使用转到槽
}

/*类的继承
*/

class Shape: public rectangle{

};


//创建socket界面

\end{minted}

% 代码写在外部
% \lstinputlisting[
%      style=cpp,
%     caption={Qt代码},
%     label={code:Qt代码}
% ]{./qt.cpp}

\section{信号与槽}
\begin{enumerate}
    \item connect函数
    \begin{minted}{cpp}
    connect(ui->,&,this,&);
    /*connect函数是连接信号和槽，并且可以连接一个信号与多个槽；
    注意这里connect只负责“连接”（建立联系），只有信号发出才会真正建立连接*/
    /*********************/
    //信号还可以连接信号
    \end{minted}
    \item 信号槽的连接方式
    \item Lambda表达式
\end{enumerate}